% !TeX root = part8.tex
\documentclass[../final.tex]{subfiles}

\begin{document}

\practice{8}{22.10.2026}{Уравнение Дирака. \texorpdfstring{$\gamma$}{γ}-матрицы}

\rtheory{Уравнение Дирака}

\begin{equation*}
    \left(i\gamma^\mu\partial_\mu-m\right)\psi=0,
    \qquad
    \partial_\mu=\left(\frac{\partial}{\partial t},\,\nabla\right).
\end{equation*}
Здесь $\psi$ --- \rconcept{спинор}, $\gamma^\mu$ --- \rconcept{матрицы Дирака}.

\subsection*{Представление Паули--Дирака}

\begin{equation*}
    \gamma^\mu=\left(\beta,\,\beta\vec\alpha\right),
    \qquad
    \vec\alpha=
    \begin{pmatrix}
        0&\vec\sigma\\
        \vec\sigma&0
    \end{pmatrix},
    \qquad
    \beta=
    \begin{pmatrix}
        I&0\\
        0&-I
    \end{pmatrix}.
\end{equation*}
\rconcept{Матрицы Паули}:
\begin{equation*}
    \sigma_1=
    \begin{pmatrix}
        0&1\\
        1&0
    \end{pmatrix},
    \qquad
    \sigma_2=
    \begin{pmatrix}
        0&-i\\
        i&0
    \end{pmatrix},
    \qquad
    \sigma_3=
    \begin{pmatrix}
        1&0\\
        0&-1
    \end{pmatrix}.
\end{equation*}

\subsection*{Свойства \texorpdfstring{$\gamma$}{γ}-матриц}

\rconcept{Антикоммутационное соотношение}:
\begin{equation*}
    \gamma^\mu\gamma^\nu+\gamma^\nu\gamma^\mu=2g^{\mu\nu}.
\end{equation*}
Компоненты:
\begin{equation*}
    \gamma^\mu=\left(\gamma^0,\gamma^1,\gamma^2,\gamma^3\right),
    \qquad
    \gamma^0=
    \begin{pmatrix}
        I&0\\
        0&-I
    \end{pmatrix},
    \qquad
    \vec\gamma=
    \begin{pmatrix}
        0&\vec\sigma\\
        -\vec\sigma&0
    \end{pmatrix}.
\end{equation*}
Квадраты:
\begin{align*}
    (\gamma^0)^2&=\mathbf 1 \quad\text{--- единичная матрица},\\
    (\gamma^k)^2&=
    \begin{pmatrix}
        0&\sigma_k\\
        -\sigma_k&0
    \end{pmatrix}
    \begin{pmatrix}
        0&\sigma_k\\
        -\sigma_k&0
    \end{pmatrix}
    =
    \begin{pmatrix}
        -\sigma_k^2&0\\
        0&-\sigma_k^2
    \end{pmatrix}
    =-\mathbf 1.
\end{align*}
Эрмитово сопряжение:
\begin{align*}
    (\gamma^0)^\dagger&=\gamma^0,\\
    (\gamma^k)^\dagger&=
    \begin{pmatrix}
        0&\sigma_k\\
        -\sigma_k&0
    \end{pmatrix}^\dagger
    =
    \begin{pmatrix}
        0&-\sigma_k^\dagger\\
        \sigma_k^\dagger&0
    \end{pmatrix}
    =-\gamma^k.
\end{align*}

\task[1]{Свёртки \texorpdfstring{$\gamma$}{γ}-матриц}
\label{task:gamma_contractions_sem8}

\condition
Обозначим $\hat a=a^\mu\gamma_\mu$ (в другой записи $a\!\!\!/$).
Упростить свёртки
\begin{equation*}
    \gamma^\mu\gamma_\mu,\qquad
    \gamma^\mu\hat a\gamma_\mu,\qquad
    \gamma^\mu\hat a\hat b\gamma_\mu,\qquad
    \gamma^\mu\hat a\hat b\hat c\gamma_\mu.
\end{equation*}

\solution
\subsection*{1) \texorpdfstring{$\gamma^\mu\gamma_\mu$}{γ^μ γ_μ}}

\begin{equation*}
    \gamma^\mu\gamma_\mu=\gamma^\mu\gamma^\nu g_{\mu\nu}
    =\frac12\left(\gamma^\mu\gamma^\nu+\gamma^\nu\gamma^\mu\right)g_{\mu\nu}
    =\frac12\cdot2g^{\mu\nu}g_{\mu\nu}=4.
\end{equation*}

\subsection*{2) \texorpdfstring{$\gamma^\mu\hat a\gamma_\mu$}{γ^μ a γ_μ}}

\begin{align*}
    \gamma^\mu\hat a\gamma_\mu
    &=a_\nu\gamma^\mu\gamma^\nu\gamma_\mu
    =a_\nu\left(2g^{\mu\nu}-\gamma^\nu\gamma^\mu\right)\gamma_\mu\\
    &=2a^\mu\gamma_\mu-a_\nu\gamma^\nu\cdot4
    =2\hat a-4\hat a=-2\hat a.
\end{align*}

\subsection*{3) \texorpdfstring{$\gamma^\mu\hat a\hat b\gamma_\mu$}{γ^μ a b γ_μ}}

\begin{align*}
    \gamma^\mu\hat a\hat b\gamma_\mu
    &=a_\nu\gamma^\mu\gamma^\nu\hat b\gamma_\mu
    =a_\nu\left(2g^{\mu\nu}-\gamma^\nu\gamma^\mu\right)\hat b\gamma_\mu\\
    &=2a^\mu\hat b\gamma_\mu-a_\nu\gamma^\nu\left(\gamma^\mu\hat b\gamma_\mu\right)
    =2\hat b\hat a+2\hat a\hat b
    =2\left(\hat b\hat a+\hat a\hat b\right)\\
    &=2\left(\gamma_\mu a^\mu\gamma_\nu b^\nu+\gamma_\nu b^\nu\gamma_\mu a^\mu\right)
    =2a^\mu b^\nu\left(\gamma_\mu\gamma_\nu+\gamma_\nu\gamma_\mu\right)\\
    &=4a^\mu b^\nu g_{\mu\nu}=4(a\cdot b).
\end{align*}

\subsection*{4) \texorpdfstring{$\gamma^\mu\hat a\hat b\hat c\gamma_\mu$}{γ^μ a b c γ_μ}}

\begin{align*}
    \gamma^\mu\hat a\hat b\hat c\gamma_\mu
    &=a_\nu\gamma^\mu\gamma^\nu\hat b\hat c\gamma_\mu
    =a_\nu\left(2g^{\mu\nu}-\gamma^\nu\gamma^\mu\right)\hat b\hat c\gamma_\mu\\
    &=2a^\mu\hat b\hat c\gamma_\mu-a_\nu\gamma^\nu
    \left(\gamma^\mu\hat b\hat c\gamma_\mu\right)
    =2\hat b\hat c\hat a-2\hat a\left(\hat c\hat b+\hat b\hat c\right)\\
    &=2\hat b\hat c\hat a-2\hat a\hat c\hat b-2\hat a\hat b\hat c.
\end{align*}
Скобка $\hat c\hat b+\hat b\hat c=2(b\cdot c)$ --- \rnote{число}, поэтому её
можно переставить через $\hat a$:
\begin{equation*}
    \gamma^\mu\hat a\hat b\hat c\gamma_\mu
    =2\hat b\hat c\hat a-2\left(\hat c\hat b+\hat b\hat c\right)\hat a
    =-2\hat c\hat b\hat a.
\end{equation*}

\answer
\begin{align*}
    \gamma^\mu\gamma_\mu&=4, &
    \gamma^\mu\hat a\gamma_\mu&=-2\hat a,\\
    \gamma^\mu\hat a\hat b\gamma_\mu&=4(a\cdot b), &
    \gamma^\mu\hat a\hat b\hat c\gamma_\mu&=-2\hat c\hat b\hat a.
\end{align*}
Для нечётного числа множителей в общем виде
\begin{equation*}
    \gamma^\mu\hat a_1\hat a_2\ldots\hat a_{2n+1}\gamma_\mu
    =-2\hat a_{2n+1}\hat a_{2n}\ldots\hat a_1.
\end{equation*}

\rtheory{Матрица \texorpdfstring{$\gamma^5$}{γ5}}

\begin{equation*}
    \gamma^5=i\gamma^0\gamma^1\gamma^2\gamma^3=
    \begin{pmatrix}
        0&I\\
        I&0
    \end{pmatrix}.
\end{equation*}
Свойства:
\begin{align*}
    (\gamma^5)^2&=
    \begin{pmatrix}
        0&I\\
        I&0
    \end{pmatrix}
    \begin{pmatrix}
        0&I\\
        I&0
    \end{pmatrix}
    =\mathbf 1,\\
    (\gamma^5)^\dagger&=\gamma^5,\\
    \gamma^5\gamma^\mu+\gamma^\mu\gamma^5&=0.
\end{align*}

\task[2]{Следы \texorpdfstring{$\gamma$}{γ}-матриц}
\label{task:gamma_traces_sem8}

\condition
Вычислить следы $\operatorname{Tr}\mathbf 1$ и $\operatorname{Tr}[\hat a\hat b]$.

\solution
\begin{enumerate}[label=\arabic*)]
    \item $\operatorname{Tr}\mathbf 1=4$.
    \item $\operatorname{Tr}[\hat a\hat b]$ вычисляется на следующем занятии
    с помощью общей формулы для следа произведения $\gamma$-матриц.
\end{enumerate}

\end{document}
